\section{Trabajo realizado}

Durante el desarrollo de la aplicación se priorizó la disponibilidad de una demostración funcional o \textit{mínimo producto viable} (MVP) para la presentación final.  
Este enfoque permitió mostrar las funcionalidades principales de la aplicación, teniendo en cuenta que se optó por implementar una idea original y que el tiempo de desarrollo disponible fue limitado en relación con el \textit{backlog} inicialmente planteado.

\subsection{Backend}

En el backend se ha logrado implementar la mayoría de las funcionalidades necesarias para el correcto funcionamiento de la aplicación, permitiendo la gestión de usuarios, la comunicación con la base de datos y la exposición de los servicios requeridos por el frontend.

\subsection{Frontend}

En el frontend también se han implementado la mayor parte de las funcionalidades previstas inicialmente. Entre las principales características desarrolladas se encuentran:

\begin{itemize}
    \item Operaciones CRUD de \textit{havens}.
    \item Visualización del estado de la cuenta del usuario y de la disponibilidad para la creación de nuevos \textit{havens}.
    \item Sistema de mensajería (chat).
    \item Subida y descarga de imágenes.
    \item Funcionalidades de geolocalización.
\end{itemize}

\section{Trabajo pendiente}

A pesar de los avances realizados, existen diversas funcionalidades que no han podido ser implementadas debido a las limitaciones de tiempo y alcance del proyecto.

\subsection{Backend}

En el backend quedan pendientes las siguientes mejoras y ampliaciones:

\begin{itemize}
    \item Implementación de un sistema de autenticación de doble factor mediante correo electrónico u otros mecanismos.
    \item Sistema de actualización de la cuenta a modo \textit{premium} basado en el resultado de una pasarela de pago, en lugar de una actualización directa en la base de datos.
    \item Integración completa de una pasarela de pago.
    \item Mejoras adicionales de seguridad.
    \item Implementación de WebSockets para la generación de notificaciones en tiempo real.
\end{itemize}

\subsection{Frontend}

En el frontend han quedado pendientes diversos aspectos relacionados con la experiencia de usuario (\textit{user experience}):

\begin{itemize}
    \item Actualización del sistema de chat en tiempo real.
    \item Uso de WebSockets para la recepción de notificaciones.
    \item Implementación de un radar de \textit{havens}. Inicialmente, los \textit{havens} cercanos no estaban pensados para mostrarse en una lista, sino mediante una visualización tipo radar. Debido a las restricciones de tiempo, esta funcionalidad no pudo ser desarrollada.
\end{itemize}

\subsection{Mejoras adicionales}
Además de las funcionalidades específicas pendientes, existen diversas mejoras generales que podrían 
implementarse para optimizar la aplicación:
\begin{itemize}
    \item Mejoras en la interfaz de usuario para una experiencia más intuitiva y atractiva.
    \item Optimización del rendimiento tanto en el backend como en el frontend.
    \item Ampliación de las pruebas unitarias y de integración para asegurar la calidad del código.
    \item Documentación más detallada del código y de la arquitectura del sistema.
    \item Implementación de un sistema de análisis y monitoreo para el backend.
    \item Implementación de despliegue continuo (CI/CD) para facilitar futuras actualizaciones.
    \item Soporte para múltiples idiomas en la aplicación móvil.
    \item Implementación de metodos de monetización adicionales, como suscripciones o anuncios.
\end{itemize}